\documentclass{article}
\usepackage{amsmath}
\begin{document}
\section*{A qualified relative reading can make independent clocks usable}
Two early-armed observers must each carry out a calibration in
the approved interval from physical tick 3 through tick 9, while
preserving an old audit at either tick 4 or tick 6. Each clock
has constant offset in \(\{1,-2\}\), independently possible.
A single common numeric target is not sufficient even if the
planner knows the actual offsets in a mixed world. For example,
in the world \((1,-2)\) the earliest bound requires an early
common target at least 4 but the first audit would require it
at most 2. A common target for the late branch would have to
exceed 7 but the deadline requires it to be at most 7. The
opposed mixed world is symmetric.

The action interface can instead control the two local thresholds
separately. A qualified clock operation samples both clocks
at one modeled physical cut and exports only
\(k=\delta_B-\delta_A\). No raw absolute offset leaves this
trusted calibration source. Returning local thresholds
\((5,5+k)\) makes the physical activations coincide at
\(5-\delta_A\), tick 4 or 7, so it already proves feasibility.
For minimum worst-case latest act, the mixed classes use 4
or 1 instead of 5 on A (and add the same shift on B), which
acts at tick 3. The equally informed direct search compares
all target vectors and gives the same result. The relative
reading can identify the pair in mixed cases of this small
source alphabet; the wire claim is not absolute-clock privacy.
All these targets retain both deadlines and the old audit. The source must own the
same-time read and the offsets must remain fixed after it;
independent remote sample timestamps with unknown transport
delay do not produce that proof.

\paragraph{Sharp finite decision information requirement.}
For this deterministic source alphabet and one-shot vector
interface, the two mixed worlds cannot share an admissible
output with each other or with either equal-offset world.
This follows by intersecting the respective early and late
target boxes in the direct finite check. Neither coordinate
needs to be searched above 10 by the deadline bound, or below
zero in this command interface. The two equal-offset worlds
do share \((5,5)\). Consequently any exact zero-error plan
selector needs at least three distinction classes. The
three values of \(k\) realize exactly three. This is not
a lower bound for other protocols that reveal on-board state,
move the protected audits or change the allowed physical action
interface.

\paragraph{Actual scoped delivery.}
The test routes the relative value through actual native
rolling credits and arms each observer with its own addressed
command only if the qualified value arrived. An independently
owed five-byte delivery leaves room for that one-byte value.
A six-byte delivery blocks it, although the newer frame
attempt uses transmit credit. The default then withholds the
new calibrations and keeps the old audit; it does not claim
both deadlines. Across all four offset pairs with both prior
credit cases, the qualified and matched direct controls give
the same physical outcomes. All four qualified cases achieve
the full service, all four blocked cases correctly preserve the
old audit under deferral, and the raw common-target control
loses the old audit on four split-clock histories.
The held-in native fixture makes no claim to hardware
synchronization, broader clock stability, independent mission
acceptance, or advantage over an equally qualified direct timer.
\end{document}
